\paragraph{Report the declared accounts and explain their movements.}
A common book also needs a reporting method that respects each company's chart.
\texttt{CompanyBookReport/1} takes opening and closing books for the same company,
period and account population, together with a declared statement map. Every
account belongs to exactly one line of the same classification: asset, liability,
equity, income or expense. A missing account, overlapping line or conflicting
classification prevents the report. Different memory schemas can describe the
opening and closing books, since a period may create a new lot; each book is
validated independently, and reporting does not establish that memory migration.

The report follows the book's debit-positive convention. Assets and expenses use
their stored signs; liabilities, direct equity and income reverse them. Profit
equals income less expenses, and reported equity includes that profit. Contra
balances remain signed. The first profile requires an opening book at the period
start with closed income and expense accounts and no current-period cash flows.
It therefore describes a recorded period, rather than an arbitrary comparison
between two accumulated income statements.

Cash flows preserve each method's declared activity and origin. For operating,
investing and financing activities separately, the report sums receipts and
payments before computing the net movement. The aggregate must equal the change
in the selected cash account. Each flow must refer to an event recorded after
the opening history; a prior event cannot be reused to explain current-period
cash. Direct equity movements require explicit account, event, date, amount and
description. They must explain every change in the direct equity accounts.
Profit enters the equity reconciliation separately, so a missing contribution
or distribution cannot disappear into an unexplained residual.

An authored sale and contribution illustrate all four statements. Opening cash
is USD~100, inventory USD~20, another asset USD~10, liabilities USD~30 and equity
USD~100. Selling inventory with carrying amount USD~15 for USD~40 produces
USD~25 profit. A subsequent USD~5 owner contribution brings cash to USD~145.
Closing assets of USD~160 equal liabilities of USD~30 plus equity of USD~130;
opening equity of USD~100, the USD~5 contribution and USD~25 profit explain that
closing equity. Operating cash receipts are USD~40 and financing receipts USD~5.
The same quantities reconcile under two distinct company charts. In a compatible
retained annual case, every asset, liability, income and expense line and all
direct cash-flow activities agree with the native report: profit is USD~49.50,
cash USD~491.50 and equity USD~671.50.

These equalities establish the recorded amounts' consistency. They cannot reveal
omitted offsetting events or establish the economic cause of a supplied equity
explanation. The report therefore retains its recorded-period scope and never
claims fiscal completeness. Note requirements remain explicitly missing,
unsupported or unresolved; this profile cannot declare a note package complete.
The indirect operating-cash bridge, source and accounting admission, comparative
presentation and explicit closing and carry methods remain separate obligations.
